% !PRINTED 266
\weberpage{266}

%<<CONT
and have learned the ways by which one can arrive at their computation in other cases as well.

We have already remarked earlier that the Bezoutiant is not a general quadratic form, but that it has the special peculiarity that the number $\nu$ of its negative squares can never become greater than $\frac{1}{2}n$. This peculiarity must find its expression in certain inequality conditions between the coefficients, which, however, are not yet known at present. Only in the case $n = 3$ can we state the algebraic character of this restriction completely.

For the cubic equation
\[
f(x) = a_0 x^3 + a_1 x^2 + a_2 x + a_3 = 0
\]
the Bezoutiant is, by \S.\,73, (13), nothing other than the \gesperrt{Hessian} covariant
\[
-\tfrac{2}{3} H(t_1, t_0) = -\tfrac{2}{3}(A_0 t_1^2 + A_1 t_1 t_0 + A_2 t_0^2).
\]
But if $D$ is the discriminant, so that
\[
3D = 4A_0 A_2 - A_1^2
\]
there holds the relation [\S.\,62, (4), (12)]
\[
4H^3 + Q^2 + 27Df^2 = 0,
\]
and hence, if we set $t_0 = 0$ in it,

\weq{1}{-4A_0^3 + q_0^2 + 27Da_0^2 = 0,}

where the coefficient appearing above is given by
\[
q_0 = 27a_0^2 a_3 - 9a_0 a_1 a_2 + 2a_1^3\text{.}
\]
This relation shows that $A_0$ and $D$ cannot both be positive at the same time; and hence that the form $-H$ is not decomposable into two negative squares.

We wish to apply the general theorem of the preceding section to the biquadratic equation as well, which for the sake of simplicity we assume in the form

\weq{2}{x^4 + ax^2 + bx + c = 0}

To form the coefficients of the Bezoutiant, we have, in the formulas \S.\,73, (14), to replace
\[
a_0,\ a_1,\ a_2,\ a_3,\ a_4
\]
by
\[
1,\ 0,\ a,\ b,\ c
\]
We thus obtain
\[
\begin{aligned}
&B_{2,2} = -2a, \quad B_{1,1} = a^2 - 4c, \quad B_{0,0} = -\tfrac{1}{4}(8ac - 3b^2), \\
&B_{0,1} = \tfrac{1}{2}ab, \quad B_{2,0} = -4c, \quad B_{1,2} = -3b.
\end{aligned}
\]
% [resolved 2026-08-06, 8x zoom strip 8/8] The smudged sign inside B_{0,0} is a MINUS.
% Compared against the clean minus signs on these same two lines (-2a, -4c, -3b, -1/4) and
% the plus signs three lines above (-4A_0^3 + q_0^2 + 27 D a_0^2): the disputed stroke has
% no vertical bar, only a horizontal one with a speck of ink. B_{0,0} = -1/4(8ac - 3b^2)
% stands as transcribed. Same strip also confirms B_{2,2} = -2a, B_{1,1} = a^2 - 4c,
% B_{0,1} = 1/2 ab, B_{2,0} = -4c, B_{1,2} = -3b.
